\documentclass{report}
\usepackage{amsfonts, amsmath}
\newcommand\R{{\mathbb R}}
\newcommand\Q{{\mathbb Q}}
\newcommand\N{{\mathbb N}}
\newcommand{\vf}{\varphi}
\pagestyle{empty}
\begin{document}
\hfuzz 12pt
\large
\centerline{\Huge \bf Analisi Matematica II}
\vskip.2cm
\centerline{\Huge Correzione Terzo Appello}

\vskip1cm
\begin{enumerate}
\item Sia $x(n)$ il capitale dopo $n$ bimestri, chiaramente
$x(0)=x$. Sia inoltre $6\Delta$ il tasso annumale praticato, dunque
ogni bimestre viene pagato il $\Delta$ per cento del capitale. Cio\`e
\[
x(n+1)=x(n)+\Delta x(n)=(1+\Delta)x(n)=(1+\Delta)^nx.
\]
Sei anni corrispondono a $n=36$ bimestri, dunque
\[
1.4 x=(1+\Delta)^{36} x.
\]
Passando ai logaritmi, 
\[
\ln (1+\Delta)=\frac {\ln 1.4}{36}.
\]
a questo punto se avete un calcolatore avete finito (e ottenete il 5.6
per cento annuale, non male), altrimenti ricordate la serie di Taylor 
\[
\ln(1+a)=a-\frac 12 a^2+\frac 13 a^3- ...
\]
Usando tale formula al secondo ordine si ha $\ln 1.4\sim .32$  (invece
di $.336...$). Usando poi la formula al primo ordine per $\ln
(1+\Delta)$, visto che $\Delta$ \`e pi\`u piccolo, si ha
$\Delta\sim 2/225$, dunqe l'interesse
percentuale annuo sarebbe stimato da $600\Delta\sim 5.3$,
che \`e essenzialmente giusto.
\item Integrando per parti si ha
\[
\int_0^1f(x)\sin\lambda x
dx=\frac{f(0)-f(1)\cos\lambda}{\lambda}+\frac
1\lambda\int_0^1f'(x)\cos\lambda x dx
\]
poich\`e $f$ \`e due volte derivabile $f'$ \`e derivabile, dunque
l'integrale \`e be definito. A questo punto, poich\`e $|\cos\lambda
x|\leq 1$ il risultato \`e ovvio.

\item Per l'argomento che si usa nel criterio integrale si ha
\[
\sum_{n=k}^\infty< \lim_{L\to\infty}\int_{k-1}^L\frac 1{x^2}
dx=\frac 1{k-1}.
\]
Dunque $\sum_{n=11}^\infty \frac 1{n^2}< \frac 1{10}$. Basta quindi
sommare i primi dieci termini per avere la risposta (per esempio $3$
\`e una risposta corretta).
\item Poich\`e in un intorno di zero $y$ non sar\`a zero si puo
dividere ed ottenere
\[
\frac{y'(x)}{y(x)}=e^x
\]
integrando poi tra zero e $x$, e ricordano le condizioni iniziali, si ha
\[
\ln y(x)=e^x-e
\]
cio\`e $y(x)=e^{e^x-e}$.
\item Poich\`e
\[
\int_0^1\frac {x+3}{x^2+1} dx=\frac
12\int_0^1\frac{2x}{x^2+1}dx+3\int_0^1\frac 1{x^2+1}dx
\]
ora nel primo integrale applichiamo la sostituzione $z=x^2$, mentre il
secondo \`e la primitiva di $\arctan$. Dunque
\[
\int_0^1\frac {x+3}{x^2+1} dx=\frac
12\int_0^1\frac 1{1+z}dz+3\arctan 1=\frac 12\ln 2+3\arctan 1.
\]
\item Il generico polinomio di primo grado ha la forma $p(x)=ax+b$,
dunque $A(p)=A(a,b)$. Supponiamo $a$ fissato e vediamo i punti criti
rispetto a $b$
\[
\frac{dA(a,b)}{db}=2\int_{-\pi}^{\pi}(\sin x-ax-b)dx=-2\pi b,
\]
dunque $b=0$ \`e il solo punto critico e poich\`e $A$ \`e sempre maggiore di
zero, deve essere di minimo. Si possono perci\`o considereare solo il
polinomi del tipo $ax$. Vediamo dunque quali $a$ danno un punto
critico
\[
\frac{dA(a,0)}{da}=2\int_{-\pi}^{\pi}(\sin x-ax)x dx=4\pi-\frac43\pi^3a
\]
(l'integrale di $x\sin x$ si fa
istantaneamente per parti). Di nuovo c'e' un solo punto critico che
quindi deve essere di minimo. Dunque $p(x)=3\pi^{-2}x$ rende  minimo $A$.
\end{enumerate}
\end{document}
